\DocumentMetadata{lang=en,tagging=on,
% pdfstandard=ua-1,pdfversion=1.7,tagging-setup={math/alt/use}
 pdfstandard=ua-2,tagging-setup={math/setup=mathml-SE}
}
\documentclass[twocolumn]{article}
\usepackage{graphicx}
\usepackage{amsmath}
\usepackage{booktabs}
\usepackage{hyperref}
\title{Sample Document}
\author{Bob Tomper}
\date{}
\begin{document}
\maketitle

\tableofcontents

\section{Some Items}
This is a sample document that has various items that can appear in an ordinary\footnote{for some definitions of ordinary} mathematics paper.  The Pythagorean Theorem says:
\begin{equation}a^2+b^2=c^2.\end{equation}
The inverse of a $2\times2$ matrix is
\begin{equation}\begin{pmatrix}a&b\\c&d\end{pmatrix}^{-1}=\frac1{ad-bc}\begin{pmatrix}d&-b\\-c&a\end{pmatrix}.\end{equation}
\subsection{Floats}
\autoref{tab:tagging} shows the tagging status of various \TeX\ engines.
\autoref{fig:duck} shows an example image.

\begin{table}[hbp]
\centering
\tagpdfsetup{table/header-rows={1}}
\caption{\label{tab:tagging}TeX engine tagging status.}
\begin{tabular}{cc}\toprule
Engine & Taggable\\\midrule
Lua\TeX & yes \\
\TeX & somewhat \\
Xe\TeX & no \\\bottomrule
\end{tabular}
\end{table}

\begin{figure}\centering
\includegraphics[width=\linewidth,alt={The word Image over a light gray background image with four gray lines intersecting.}]
 {example-image}
\caption{\label{fig:duck}An example image}
\end{figure}

\section{Citations}
Here we cite \cite{Knuth:TB5-1-67} and \cite{Hoekwater:TB28-3-312} from the \texttt{tugboat.bib} sample file.

\bibliographystyle{plain}
\bibliography{tugboat}

\end{document}
